\documentclass[10pt,a4paper]{amsart}
\usepackage[margin=19mm]{geometry}
\usepackage{amsmath,amssymb,mathtools}
\usepackage[scr=rsfs,cal=euler]{mathalfa}
\usepackage{xfrac}
\usepackage[unicode,hidelinks,bookmarks=false]{hyperref}
\newcommand{\ie}{i.e.\ }
\newcommand{\cf}{cf.\ }
\newcommand{\resp}{respectively\ }
\newcommand{\Subsection}{\subsection}
\providecommand{\nobreakdash}{\nobreak\hbox{-}}
\setlength{\emergencystretch}{2em}
\multlinegap0pt
\usepackage{subfig}
\usepackage{tikz}
\usepackage{float}
\usepackage{multirow}
\usepackage{algorithm}\usepackage{algpseudocode}
\newtheorem{theorem}{Theorem}
\title{Accelerated Online Risk-Averse Policy Evaluation in POMDPs with Theoretical Guarantees and Novel CVaR Bounds}
\author{Yaacov Pariente, Vadim Indelman}
\date{Source: arXiv:2602.23073v1 (CC BY 4.0)}
\begin{document}
\maketitle

\noindent\emph{Source and licence.} Accelerated Online Risk-Averse Policy Evaluation in POMDPs with Theoretical Guarantees and Novel CVaR Bounds. Version arXiv:2602.23073v1 (CC BY 4.0), distributed by arXiv under the Creative Commons Attribution 4.0 International licence, http://creativecommons.org/licenses/by/4.0/, as recorded in the arXiv licence metadata for this exact version and shown on the abstract page https://arxiv.org/abs/2602.23073v1. Full licence notice: \texttt{licenses/arxiv-2602.23073-CC-BY-4.0.txt}.\par\smallskip

\noindent\textit{Editorial scope.} Counted unit: the article's theorem prf:simplification\_bound\_over\_belief\_distribution (source0.tex lines 1880--1895). The article's complete original proof of it is delivered with it (source0.tex lines 1896--2024, one \texttt{proof} environment of 129 lines). No mathematics is written, repaired or re-derived: the statement and the proof are the article's own lines, in the article's own notation, and no retained line is substituted or re-flowed.  No further source-local context is needed: every cross-reference of every retained line resolves inside the two delivered spans.  Nothing was omitted from any retained span, so the package carries no omission record.  No retained line contains a citation, so the wrapper carries no bibliography and no bibliography entry.  No retained line contains a figure environment, an image-inclusion command or drawing code, so no image is delivered and the package declares no asset dependency.  Editorial formatting only, in addition: an AMS \texttt{amsart} preamble that restates the article's own declarations and macros used by the retained lines, namely the environments \texttt{theorem} on separate counters, together with the standard packages those lines need.  The retained environments print in this document as Theorem 1 in order of appearance, whereas the article prints the same items with the numbering of its own full document.  The complete native source of the article as distributed in the arXiv e-print for this exact version is bundled unchanged as \texttt{sources/195369/source0.tex}.
\begin{theorem}\label{prf:simplification_bound_over_belief_distribution}
	\begin{equation}
		\begin{aligned}
			&\sup_{l\in \mathbb{R}}|P(R_{k:T} \leq l|b_k, a_k,\pi) \! -\! P_s(R_{k:T} \leq l|b_k, a_k,\pi)| \!
			\leq \!\! \! \sum_{t=k}^{T-1} \mathbb{E}^s[\Delta^s(b_t, a_t)|b_k, a_k],
		\end{aligned}
	\end{equation}
	where $\Delta^s$ is the TV distance that is defined by
	\begin{equation}
		\begin{aligned}
			&\Delta^s(b_{t-1},a_{t-1}) \triangleq 
			\int_{b_t\in B} |P(b_t|b_{t-1},a_{t-1})-P_s(b_t|b_{t-1},a_{t-1})|db_t.
		\end{aligned}
	\end{equation}
	
\end{theorem}

\begin{proof}
	By expanding the belief path from time $k+1$ to time $T$ we get
	\begin{equation}
		\begin{aligned}
			P(\sum_{t=k}^T c(b_t,a_t)\leq l|b_{k},\pi) 
			& = \int_{b_{k+1:T}\in B^{T-k}} P(R_{k:T}\leq l|b_{k:T},\pi) 
			 \prod_{i=k+1}^T P(b_i|b_{i-1}, \pi, a_i)db_{k+1:T}\\
			& = \int_{b_{k+1:T}\in B^{T-k}} 1_{R_{k:T}\leq l}\prod_{i=k+1}^T P(b_i|b_{i-1}, \pi, a_i)db_{k+1:T},
		\end{aligned}
	\end{equation}
	where the second equality holds because $R_{k:T}$ is constant given $\pi$ and $b_{k:T}$. Hence, the subtraction between the simplified and original return CDFs can be exhibited using the difference between the multiplication of the belief transition models.
	\begin{equation}
		\begin{aligned}
			&P(R_{k:T}\leq l|b_k,\pi)-P_s(R_{k:T}\leq l|b_k,\pi)
			=\int_{b_{k+1:T}\in B^{T-k}} 1_{R_{k:T}\leq l} [\prod_{i=k+1}^T P(b_i|b_{i-1}, \pi, a_i) 
			- \prod_{i=k+1}^T P_s(b_i|b_{i-1}, \pi, a_i)]db_{k+1:T}
		\end{aligned}
	\end{equation} 
	By applying the triangle inequality we get the following bound
	\begin{equation}\label{eq:prf_first_belief_cdf_bound}
		\begin{aligned}
			&|P(R_{k:T}\leq l|b_k,\pi)-P_s(R_{k:T}\leq l|b_k,\pi)|
			\leq \int_{b_{k+1:T}\in B^{T-k}} 1_{R_{k:T}\leq l} |\prod_{i=k+1}^T P(b_i|b_{i-1}, \pi, a_i) 
			- \prod_{i=k+1}^T P_s(b_i|b_{i-1}, \pi, a_i)|db_{k+1:T} \\
			&\leq \int_{b_{k+1:T}\in B^{T-k}} |\prod_{i=k+1}^T P(b_i|b_{i-1}, \pi, a_i) 
			- \prod_{i=k+1}^T P_s(b_i|b_{i-1}, \pi, a_i)|db_{k+1:T} = g(T),
		\end{aligned}
	\end{equation}
	where 
	\begin{equation}
		\begin{aligned}
			g(t) &\triangleq \int_{b_{k+1:t}\in B^{t-k+1}} |\prod_{i=k+1}^t P(b_i|b_{i-1}, \pi, a_i) 
			- \prod_{i=k+1}^t P_s(b_i|b_{i-1}, \pi, a_i)|db_{k+1:t}.
		\end{aligned}
	\end{equation}
	We will prove that
	\begin{equation}
		g(t+1) \leq g(t) + \mathbb{E}^s[\Delta^s (b_{t}, a_t)|b_k, a_k]=\epsilon_t
	\end{equation}
	for $t \in \{k+1,\dots, T\}$, where $\epsilon_k=0$, and therefore 
	\begin{equation}
		|P(R_{k:t}\leq l|b_k,\pi)-P_s(R_{k:t}\leq l|b_k,\pi)|\leq \epsilon_t.
	\end{equation}
	For the case where $t=k$, $b_t$ is constant, and therefore the $c(b_t, a_t)$ does not depend on the belief transition model. Hence, $P_s(c(b_t, a_t)\leq l|b_t, a_t)=P(c(b_t, a_t)\leq l|b_t, a_t)$. Therefore we get
	\begin{equation}
		P(c(b_k, a_k)\leq l|b_k, a_k) - P_s(c(b_k, a_k)\leq l|b_k, a_k)|=0=\epsilon_k.
	\end{equation}
	The rest of the proof will use induction over $t\in \{k+1, \dots, T\}$.
	
	\textbf{Base case, $t=k+1$:}
	\begin{equation}
		\begin{aligned}
			&g(k+1)=\int_{b_{k+1}\in B} |\prod_{i=k+1}^{k+1} P(b_i|b_{i-1}, \pi, a_i) 
			- \prod_{i=k+1}^{k+1} P_s(b_i|b_{i-1}, \pi, a_i)|db_{k+1:k+1} \\
			&= \int_{b_{k+1}\in B} |P(b_{k+1}|b_{k}, \pi, a_{k}) - P_s(b_{k+1}|b_{k}, \pi, a_k)|db_{k+1} 
			= \Delta^s(b_{k}, a_k)=\mathbb{E}^s[\Delta^s(b_{k}, a_k)|b_k,a_k] \\
			&= \epsilon_{k+1}
		\end{aligned}
	\end{equation} 
	
	\textbf{Induction step:} Assume that the claim is true for $t\in \{k+1, \dots, T-1\}$ and prove for $t+1$.
	
	Our proof's strategy is to express $g(t+1)$ recursively, using $g(t)$, and separate the bound at time $t+1$ into 2 components - one at time $t$ and one at time $t+1$. We then bound the term at time $t+1$, and use the induction step to bound the recursive part of $g$ at time $t$. 
	
	For all $x_i,y_i\in \mathbb{R}$,
	\begin{equation}
		\begin{aligned}
			&\prod_{i=k+1}^{t+1} x_i-\prod_{i=k+1}^{t+1}y_i
			=\prod_{i=k+1}^{t+1} x_i - x_{t+1} \prod_{i=k+1}^t y_i + x_{t+1} \prod_{i=k+1}^t y_i - \prod_{i=k+1}^{t+1}y_i 
			=x_{t+1}(\prod_{i=k+1}^t x_i - \prod_{i=k+1}^t y_i) + (x_{t+1} - y_{t+1}) \prod_{i=k+1}^{t} y_i.
		\end{aligned}
	\end{equation}
	By denoting $x_i=P(b_i|b_{i-1},a_{i-1})$ and $y_i=P_s(b_i|b_{i-1},a_{i-1})$ we get
	\begin{equation}
		\begin{aligned}
			&\prod_{i=k+1}^{t+1} P(b_i|b_{i-1}, \pi, a_i) - \prod_{i=k+1}^{t+1} P_s(b_i|b_{i-1}, \pi, a_i)\\
			&= [P(b_{t+1}|b_{t},a_{t}) - P_s(b_{t+1}|b_{t},a_{t})] \prod_{i=k+1}^{t}P_s(b_i|b_{i-1},a_{i-1}) 
			+ P(b_{t+1}|b_{t},a_{t})[\prod_{i=k+1}^{t} P(b_i|b_{i-1}, \pi, a_i) 
			- \prod_{i=k+1}^{t} P_s(b_i|b_{i-1}, \pi, a_i)].
		\end{aligned}
	\end{equation}
	By combining the equation above with the triangle inequality, we get 
	\begin{equation}
		g(t+1)\leq A_1 + A_2
	\end{equation}
	for 
	\begin{equation}
		\begin{aligned}
			A_1 &\triangleq \int\limits_{b_{k+1:t+1}\in B^{t-k}} |P(b_{t+1}|b_{t},a_{t}) - P_s(b_{t+1}|b_{t},a_{t})| 
			\prod_{i=k+1}^{t}P_s(b_i|b_{i-1},a_{i-1})b_{k+1:t},
		\end{aligned}
	\end{equation}
	\begin{equation}
		\begin{aligned}
			A_2 &\triangleq \int\limits_{b_{k+1:t+1}\in B^{t-k}} P(b_{t+1}|b_{t},a_{t}) |\prod_{i=k+1}^t P(b_i|b_{i-1}, \pi, a_i) 
			- \prod_{i=k+1}^t P_s(b_i|b_{i-1}, \pi, a_i)|db_{k+1:t} 
		\end{aligned}
	\end{equation}
	$A_1$ is the expectation over the TV-distance of the belief at time t.
	\begin{equation}
		\begin{aligned}
			&A_1 = \mathbb{E}^s[\Delta^s (b_{t}, a_t)|b_k, a_k].
		\end{aligned}
	\end{equation}
	As we see below, $A_2$ does not depend on the belief integration at time $t+1$, and therefore equals to $g(t)$. From the induction assumption, we get that $g(t)\leq \epsilon_t$.
	\begin{equation}
		\begin{aligned}
			&A_2 = \int\limits_{b_{k+1:t}\in B^{t-k}} P(b_{t+1}|b_{t},a_{t})  |\prod_{i=k+1}^t P(b_i|b_{i-1}, \pi, a_i) 
			- \prod_{i=k+1}^t P_s(b_i|b_{i-1}, \pi, a_i)|db_{k:t+1} \\
			&= \int\limits_{b_{k+1:t}\in B^{t-k}} |\prod_{i=k+1}^t P(b_i|b_{i-1}, \pi, a_i) - \prod_{i=k+1}^t P_s(b_i|b_{i-1}, \pi, a_i)| 
			\underbrace{\int_{b_{t+1}}P(b_{t+1}|b_{t},a_{t})db_{t+1}}_{=1} db_{k+1:t} \\
			&=g(t)\leq \epsilon_{t}
		\end{aligned}
	\end{equation}
	By combining the computations of $A_1$ and $A_2$,
	\begin{equation}
		\begin{aligned}
			g(t+1)&\leq A_1 + A_2 = g(t) + \mathbb{E}^s[\Delta^s (b_{t}, a_t)|b_k, a_k].
		\end{aligned}
	\end{equation}
	The last equation completes the induction proof. Utilizing the recursive relation of $\epsilon_t$, we get the bound we want to prove.
	\begin{equation}
		\begin{aligned}
			&|P(R_{k:T} \leq l|b_k,\pi)-P_s(R_{k:T} \leq l|b_k,\pi)|\leq g(T)
			\leq g(T-1) + \mathbb{E}^s[\Delta^s (b_{T-1}, a_{T-1})|b_k, a_k] 
			\leq \sum_{t=k}^{T-1} \mathbb{E}^s[\Delta^s(b_t, a_t)|b_k, a_k].
		\end{aligned}
	\end{equation}
\end{proof}

\end{document}
